\setcounter{ex}{0}
\setcounter{paragraph}{0}
\PhanI
%\loai{Lí thuyết vận tốc}
\begin{ex}%[3P1Y2-1][Loai1]
Trong một dao động cơ điều hòa, những đại lượng nào sau đây có giá trị \textbf{không} thay đổi?
\choice
{\True Biên độ và tần số}
{Gia tốc và li độ}
{Gia tốc và tần số}
{Biên độ và li độ}
\loigiai{
Trong một dao động cơ điều hòa, biên độ và tần số có giá trị không thay đổi.
}
\end{ex} \haidong{15}
\begin{ex}%[3P1B2-1][Loai1]
Vectơ vận tốc của vật dao động điều hòa luôn
\choice
{ngược hướng chuyển động}
{\True cùng hướng chuyển động}
{hướng về vị trí cân bằng}
{hướng ra xa vị trí cân bằng}
\loigiai{
Vectơ vận tốc của vật dao động điều hòa luôn cùng hướng chuyển động.
}
\end{ex} \haidong{15}
%\loai{Biểu thức vận tốc cực đại và gia tốc cực đại} 

\begin{ex}%[3P1Y2-1][Loai1]
Trong quá trình dao động, vận tốc của vật có giá trị lớn nhất (cực đại) khi vật
\choice
{đi qua vị trí cân bằng}
{\True đi qua vị trí cân bằng theo chiều dương}
{đi qua vị trí cân bằng theo chiều âm}
{ở biên}
\loigiai{
Trong quá trình dao động, vận tốc của vật có giá trị lớn nhất (cực đại) khi vật đi qua vị trí cân bằng theo chiều dương.
}
\end{ex} \haidong{15}
\begin{ex}%[3P1Y2-1][Loai1]
Trong quá trình dao động, vận tốc của vật có giá trị bằng không (vật dừng lại tức thời) khi vật ở
\choice
{biên dương ($x = A$)}
{biên âm ($x = -A$)}
{đi qua vị trí cân bằng theo chiều âm}
{\True biên dương hoặc biên âm}
\loigiai{
Trong quá trình dao động, vận tốc của vật có giá trị bằng không (vật dừng lại tức thời) khi vật ở biên dương hoặc biên âm.
}
\end{ex} \haidong{15}

\begin{ex}%[3P1Y2-1][Loai1]
Trong quá trình dao động, vật có tốc độ cực đại khi vật
\choice
{\True đi qua vị trí cân bằng}
{đi qua vị trí cân bằng theo chiều dương}
{biên âm ($x = -A$)}
{biên dương ($x=A$)}
\loigiai{
Trong quá trình dao động, vật có tốc độ cực đại khi vật đi qua vị trí cân bằng.
}
\end{ex} \haidong{15}

\begin{ex}%[3P1Y2-1][Loai1]
Vận tốc của một chất điểm dao động điều hoà biến thiên
\choice
{cùng tần số và ngược pha với li độ}
{\True cùng tần số và vuông pha với gia tốc}
{khác tần số và vuông pha với li độ}
{cùng tần số và cùng pha với li độ}
\loigiai{
Vận tốc của một chất điểm dao động điều hoà biến thiên cùng tần số và vuông pha với gia tốc.
}
\end{ex} \haidong{15}

%\loai{Lí thuyết gia tốc}
\begin{ex}%[3P1Y2-1][Loai1]
Gia tốc của một chất điểm dao động điều hoà biến thiên
\choice
{\True cùng tần số và ngược pha với li độ}
{khác tần số và ngược pha với li độ}
{khác tần số và cùng pha với li độ}
{cùng tần số và cùng pha với li độ}
\loigiai{
Gia tốc của một chất điểm dao động điều hoà biến thiên cùng tần số và ngược pha với li độ.
}
\end{ex} \haidong{15}
\begin{ex}%[3P1B2-1][Loai1]
Một vật dao động điều hòa trên trục Ox quanh vị trí cân bằng O. Vectơ gia tốc của vật
\choice
{có độ lớn tỉ lệ nghịch với độ lớn vận tốc của vật}
{luôn hướng ngược chiều chuyển động của vật}
{luôn hướng theo chiều chuyển động của vật}
{\True có độ lớn tỉ lệ thuận với độ lớn li độ của vật}
\loigiai{
Vectơ gia tốc của vật có độ lớn tỉ lệ thuận với độ lớn li độ của vật.
}
\end{ex} \haidong{15}
\begin{ex}%[3P1Y2-1][Loai1]
Trong quá trình dao động, gia tốc của vật có giá trị cực tiểu ($- \omega^2A$) khi vật
\choice
{đi qua vị trí cân bằng}
{ở biên (dương hoặc âm)}
{ở biên âm ($x = - A$)}
{\True ở biên dương ($x = A$)}
\loigiai{
Trong quá trình dao động, gia tốc của vật có giá trị cực tiểu ($- \omega^2A$) khi vật ở biên dương ($x = A$).
}
\end{ex} \haidong{15}
\begin{ex}%[3P1Y2-1][Loai1]
Trong quá trình dao động, gia tốc của vật có giá trị cực đại ($\omega^2A$) khi vật
\choice
{đi qua vị trí cân bằng}
{ở biên (dương hoặc âm)}
{\True ở biên âm ($x = - A$)}
{ở biên dương ($x = A$)}
\loigiai{
Trong quá trình dao động, gia tốc của vật có giá trị cực đại ($\omega^2A$) khi vật ở biên âm ($x = - A$).
}

\end{ex} \haidong{15}\begin{ex}%[3P1Y2-1][Loai1]
Trong quá trình dao động, gia tốc của vật có độ lớn cực tiểu khi vật
\choice
{\True đi qua vị trí cân bằng}
{ở biên (dương hoặc âm)}
{ở biên âm ($x = - A$)}
{ở biên dương ($x = A$)}
\loigiai{
Trong quá trình dao động, gia tốc của vật có độ lớn cực tiểu (0) khi vật đi qua vị trí cân bằng.
}
\end{ex} \haidong{15}
\begin{ex}%[3P1Y2-1][Loai1]
Trong quá trình dao động, gia tốc của vật có độ lớn cực đại ($\omega^2A$) khi vật
\choice
{đi qua vị trí cân bằng}
{\True ở biên (dương hoặc âm)}
{ở biên âm ($x = - A$)}
{ở biên dương ($x = A$)}
\loigiai{
Trong quá trình dao động, gia tốc của vật có độ lớn cực đại ($\omega^2A$) khi vật ở biên (dương hoặc âm).
}
\end{ex} \haidong{15}








%\loai{Tính chất chuyển động}
\begin{ex}%[3P1B2-1][Loai1]
Khi một vật dao động điều hòa, chuyển động của vật từ vị trí cân bằng ra vị trí biên là chuyển động
\choice
{nhanh dần đều}
{chậm dần đều}
{nhanh dần}
{\True chậm dần}
\loigiai{
Khi một vật dao động điều hòa, chuyển động của vật từ vị trí cân bằng ra vị trí biên là chuyển động chậm dần.
}
\end{ex} \haidong{15}
\begin{ex}%[3P1B2-1][Loai1]
Một vật dao động điều hoà trên trục Ox, tại thời điểm nào đó vận tốc và gia tốc của vật có giá trị dương. Trạng thái dao động của vật khi đó là
\choice
{\True nhanh dần theo chiều dương}
{chậm dần đều theo chiều dương}
{nhanh dần theo chiều âm}
{chậm dần theo chiều dương}
\loigiai{
$v > 0$: vật phải đi theo chiều dương, $a > 0$: vật phải có li độ âm ($x < 0$)
$\Rightarrow$ vật đang chuyển động theo chiều dương từ biên âm ($x = -A$) về vị trí cân bằng.
}
\end{ex} \haidong{15}
%\loai{Sự thay đổi của vận tốc và gia tốc}
\begin{ex}
Phát biểu nào sau đây \textbf{sai} khi nói về dao động điều hoà?
\choice
{Gia tốc sớm pha $\pi$ so với li độ}
{\True Vận tốc và gia tốc luôn ngược pha nhau}
{Vận tốc luôn trễ pha $\dfrac{\pi}{2}$ so với gia tốc}
{Vận tốc luôn sớm pha $\dfrac{\pi}{2}$ so với li độ}
\end{ex} \haidong{15}
\begin{ex}%[3P1B2-1][Loai1]
Một chất điểm đang dao động điều hòa thì
\choice
{\True ở vị trí cân bằng, chất điểm có tốc độ cực đại và gia tốc bằng không}
{ở vị trí biên, chất điểm có tốc độ cực đại và gia tốc có giá trị đạt cực đại}
{ở vị trí cân bằng, chất điểm có vận tốc bằng không và gia tốc có giá trị đạt cực đại}
{ở vị trí biên, chất điểm có vận tốc bằng không và gia tốc bằng không}
\loigiai{
Trong dao động điều hòa: Ở vị trí cân bằng, chất điểm có tốc độ cực đại và gia tốc bằng không.
}
\end{ex} \haidong{15}
 
\begin{ex}%[3P1B2-1][Loai1]
Một vật dao động điều hòa khi đang chuyển động từ vị trí cân bằng đến vị trí biên âm thì
\choice
{\True vectơ vận tốc ngược chiều với vectơ gia tốc}
{độ lớn vận tốc và gia tốc cùng tăng}
{vận tốc và gia tốc cùng có giá trị âm}
{độ lớn vận tốc và độ lớn gia tốc cùng giảm}
\loigiai{
Khi đang chuyển động từ vị trí cân bằng đến vị trí biên âm thì vectơ vận tốc ngược chiều với vectơ gia tốc.
}
\end{ex} \haidong{15}
\begin{ex}
\immini[thm]{Một vật đang thực hiện một dao động điều hoà quanh vị trí cân bằng O. Hai vị trí biên là M, N như hình vẽ.
Trong giai đoạn nào sau đây của chuyển động thì vận tốc và gia tốc cùng chiều nhau?}{\begin{tikzpicture}[draw=Mapcolor,scale=.8]
\tkzInit[ymin=-0.75,ymax=.75,xmin=-0.5,xmax=6.5]\tkzClip
\tkzDefPoints{0/0/m, 3/0/o, 6/0/n}
\tkzDrawSegments[](m,n)
\tkzDrawPoints[size=3,fill=white](m,n,o)
\tkzLabelPoint[below](m){$M$}
\tkzLabelPoint[below](n){$N$}
\tkzLabelPoint[below](o){$O$}
\end{tikzpicture}} 
\choice
{Từ O đến M}
{\True Từ N đến O}
{Từ O đến N}
{Từ M đến N}
\end{ex} \haidong{15}

%\loai{Lí thuyết khác}
\setcounter{ex}{0}
\PhanII
\begin{ex}
Trong quá trình một vật gắn với lò xo dao động điều hoà trên trục $Ox$.
\choiceTF[t]
{\True Ở vị trí cân bằng, chất điểm có gia tốc bằng không}
{Ở vị trí cân bằng, chất điểm có vận tốc cực đại}
{\True Ở vị trí biên, chất điểm có vận tốc bằng không}
{Ở vị trí biên dương, chất điểm có gia tốc cực đại}
\loigiai{
\begin{itemchoice}
\itemch
\itemch
\itemch
\itemch
\end{itemchoice}
}
\end{ex} \haidong{15}
\begin{ex}
Trong thí nghiệm con lắc lò xo dao động điều hoà trên trục $Ox$, vật nhỏ chuyển động qua lại quanh vị trí cân bằng $O$.
\choiceTF[t]
{Vectơ vận tốc của vật luôn hướng về vị trí cân bằng}
{\True Khi qua vị trí cân bằng thì độ lớn vận tốc của vật đạt cực đại}
{Khi vật ở vị trí biên gia tốc của vật bằng không}
{\True Vectơ gia tốc của vật luôn hướng về vị trí cân bằng}
\loigiai{
\begin{itemchoice}
\itemch
\itemch
\itemch
\itemch
\end{itemchoice}
}
\end{ex} \haidong{15}
\begin{ex}
Trong dao động điều hoà của một vật gắn với lò xo trên trục $Ox$, vectơ vận tốc $\vec{v}$ và vectơ gia tốc $\vec{a}$ luôn biến thiên theo thời gian.
\choiceTF[t]
{Vectơ vận tốc $\vec{v}$ và vectơ gia tốc $\vec{a}$ là vectơ hằng số}
{\True Vectơ gia tốc $\vec{a}$ đổi chiều khi chất điểm qua vị trí cân bằng}
{Trong một chu kì không có thời điểm nào vectơ vận tốc $\vec{v}$ và vectơ gia tốc $\vec{a}$ cùng chiều với nhau}
{\True Vectơ vận tốc $\vec{v}$ hướng cùng chiều chuyển động của chất điểm}
\loigiai{
\begin{itemchoice}
\itemch
\itemch
\itemch
\itemch
\end{itemchoice}
}
\end{ex} \haidong{15}
\begin{ex}
Quan sát chuyển động của một vật nhỏ gắn vào lò xo dao động điều hoà trên trục $Ox$; xét mối tương quan hướng giữa vectơ vận tốc $\vec v$ và vectơ gia tốc $\vec a$ của vật.
\choiceTF[t]
{Vectơ vận tốc $\vec v$ và vectơ gia tốc $\vec a$ cùng chiều khi vật chuyển động ra xa vị trí cân bằng}
{Vectơ gia tốc $\vec a$ đổi chiều khi vật có li độ cực đại}
{Vectơ gia tốc $\vec a$ luôn hướng ra xa vị trí cân bằng}
{\True Vectơ vận tốc $\vec v$ và vectơ gia tốc $\vec a$ cùng chiều khi vật chuyển động về vị trí cân bằng}
\loigiai{
\begin{itemchoice}
\itemch Sai: Khi vật chuyển động ra xa vị trí cân bằng, $\vec v$ hướng ra xa còn $\vec a$ luôn hướng về cân bằng nên hai vectơ ngược chiều
\itemch Sai: Vectơ gia tốc đổi chiều khi vật đi qua vị trí cân bằng, không phải tại li độ cực đại
\itemch Sai: Trong dao động điều hoà, vectơ gia tốc $\vec a$ luôn hướng về vị trí cân bằng
\itemch Đúng: Khi vật chuyển động về vị trí cân bằng, cả $\vec v$ lẫn $\vec a$ đều hướng về cân bằng nên chúng cùng chiều
\end{itemchoice}
}
\end{ex}

 

